\clearpage
\fancyhead[L]{\small NOVA ACTA UPSALIENSIS.}
\hypertarget{native-609-lower}{}
\begin{center}\small Göttingische gelehrte Anzeigen. No. 70. Pages 692..694. 2 May 1818.\\[3pt]\rule{.72\linewidth}{.3pt}\end{center}

\emph{Nova Acta regiae societatis scientiarum Upsaliensis.} Vol. VII. Upsala. 394 pages in quarto.

\begingroup\brokenpenalty=10000\relax
This collection opens with a paper left by the late Prof. \textsc{Erich Prosperin} on the orbits that satellites would describe around the Sun if their primary planets were suddenly destroyed or ceased to act on them.
Depending on the circumstances, these orbits can become ellipses, parabolas, or hyperbolas, either direct or retrograde.
The type of conic section is determined partly by the satellite's distance from its primary planet and partly by where the satellite is in its orbit at the moment of the primary planet's assumed destruction.
If, with the author, we neglect the ellipticity of the primary planet's orbit around the Sun and that of the satellite's original orbit around its primary planet, and set the product of the primary planet's distance from the Sun and its mass (taking the Sun's mass as unity) $= d$ and the distance of the satellite from its primary planet $= D$, then after the destruction of the primary planet the satellite's new orbit around the Sun will always be an ellipse if $D$ is greater than $(3+\BareRoot{8})d$.
It will always become a hyperbola if $D$ is less than $(3-\BareRoot{8})d$.
Between the two limits it can become an ellipse, a parabola, or a hyperbola.
If $D$ is greater than $d$, the new orbit cannot become retrograde.
If $D$ is greater than $d\BareRoot{\tfrac{1}{2}}$, it cannot, at any rate, become a retrograde hyperbola.
Finally, if $D$ is less than $d\BareRoot{\tfrac{1}{8}}$, it cannot remain a direct ellipse.
These propositions assume that the mass of the primary planet is very small compared with that of the Sun and that the original motion of the satellite takes place in the plane of the primary planet. Without these assumptions, the propositions require some modifications.
Mr. \textsc{Prosperin} has considered each individual satellite of our solar system separately and has taken the trouble to calculate very extensive tables for them.
The orbit of our Moon would remain a direct ellipse if it were suddenly to lose the Earth, wherever it happened to be. The satellites of Jupiter, Saturn, and Uranus, on the other hand, could also acquire parabolic or hyperbolic orbits, depending on the circumstances.
Clearly, these propositions can now also be reversed. \SourceLeaf{610}A celestial body that came into the vicinity of Jupiter, Saturn, or Uranus on a parabolic or hyperbolic orbit would enter upon a circular orbit around the planet if the planet's action began \emph{all at once}, only after the body had already drawn near.
Since that is not what happens in nature, the whole investigation amounts to nothing more than a mathematical plaything.
The author did not altogether overlook that circumstance, but he did not give it enough consideration, for it is just as \emph{impossible} for a comet that happens to come into the vicinity of a primary planet to become its satellite through the forces of attraction alone as it is for an existing satellite to leave its primary planet and depart on a parabolic or hyperbolic orbit.
\par\clearpage\endgroup
